% =====================================================================
%  Préambule : langue, mise en page, styles. Rien à modifier ici
%  pour un usage normal.
% =====================================================================

% ---- Moteur : pdfLaTeX ou XeLaTeX/LuaLaTeX --------------------------
\usepackage{iftex}
\ifPDFTeX
  \usepackage[T1]{fontenc}
  \usepackage[utf8]{inputenc}
  \usepackage{lmodern}
  \usepackage[english,french]{babel}
\else
  \usepackage{fontspec}
  \usepackage{polyglossia}
  \setmainlanguage{french}
  \setotherlanguage{english}
  \providecommand{\og}{\guillemotleft}
  \providecommand{\fg}{\unskip\guillemotright}
\fi
\usepackage{textcomp}

% ---- Mise en page (A4, prête à imprimer et relier) ------------------
\usepackage[a4paper,top=2.2cm,bottom=2.2cm,left=2.5cm,right=2.2cm,headheight=15pt]{geometry}
\usepackage{setspace}
\setstretch{1.15}
\setlength{\parskip}{0.3em}
\usepackage{microtype}
\usepackage{emptypage}

% ---- Couleurs, figures, tableaux -------------------------------------
\usepackage[table]{xcolor}
\definecolor{bleuSentinel}{HTML}{1F4E79}
\definecolor{grisClair}{HTML}{F2F4F7}
\definecolor{grisTexte}{HTML}{5A6270}
\definecolor{rougeAfaire}{HTML}{C62828}
\usepackage{graphicx}
\usepackage{float}
\usepackage[font=small,labelfont=bf]{caption}
\usepackage{subcaption}
\usepackage{booktabs}
\usepackage{tabularx}
\usepackage{longtable}
\usepackage{array}
\usepackage{multirow}
\usepackage{makecell}
\renewcommand{\arraystretch}{1.25}
\newcolumntype{L}[1]{>{\raggedright\arraybackslash}p{#1}}
\newcolumntype{Y}{>{\raggedright\arraybackslash}X}
\usepackage{enumitem}
\setlist{itemsep=0.15em, topsep=0.3em}
\usepackage{amsmath,amssymb}
\usepackage{tikz}
\usetikzlibrary{shapes.geometric,arrows.meta,positioning,fit,backgrounds,calc}
\usepackage{pgfgantt}
\usepackage[most]{tcolorbox}
\usepackage{eso-pic}
\usepackage{lastpage}

% ---- Code source ----------------------------------------------------
\usepackage{listings}
\lstset{
  basicstyle=\ttfamily\footnotesize,
  backgroundcolor=\color{grisClair},
  frame=single, rulecolor=\color{grisClair!70!black},
  numbers=left, numberstyle=\tiny\color{grisTexte}, numbersep=6pt,
  breaklines=true, breakatwhitespace=true, columns=fullflexible,
  keywordstyle=\color{bleuSentinel}\bfseries,
  commentstyle=\color{grisTexte}\itshape,
  stringstyle=\color{teal!70!black},
  showstringspaces=false, tabsize=2, xleftmargin=1.2em, framexleftmargin=1em,
  captionpos=b, aboveskip=0.8em, belowskip=0.8em
}
\renewcommand{\lstlistingname}{Extrait}
\renewcommand{\lstlistlistingname}{Liste des extraits de code}
\lstdefinelanguage{yaml}{morekeywords={true,false}, sensitive=false, morecomment=[l]{\#}, morestring=[b]"}
\lstdefinelanguage{JavaScript}{morekeywords={const,let,var,function,return,if,else,for,await,async,new,export,import,from,throw,try,catch},
  sensitive=true, morecomment=[l]{//}, morecomment=[s]{/*}{*/}, morestring=[b]', morestring=[b]"}

% ---- En-têtes et pieds de page ----------------------------------------
\usepackage{fancyhdr}
\pagestyle{fancy}
\fancyhf{}
\fancyhead[L]{\small\color{grisTexte}\nouppercase{\leftmark}}
\fancyhead[R]{\small\color{grisTexte}\titreCourt}
\fancyfoot[C]{\small\thepage}
\renewcommand{\headrulewidth}{0.4pt}
\fancypagestyle{plain}{\fancyhf{}\fancyfoot[C]{\small\thepage}\renewcommand{\headrulewidth}{0pt}}

% ---- Titres de chapitres --------------------------------------------
\usepackage{titlesec}
\titleformat{\chapter}[display]
  {\normalfont\bfseries\color{bleuSentinel}}
  {\filleft\Large\MakeUppercase{\chaptertitlename}~\thechapter}
  {1ex}
  {\titlerule[1.2pt]\vspace{1.2ex}\filright\huge}
  [\vspace{0.5ex}\titlerule]
\titleformat{name=\chapter,numberless}[display]
  {\normalfont\bfseries\color{bleuSentinel}}{}{0pt}
  {\titlerule[1.2pt]\vspace{1.2ex}\filright\huge}[\vspace{0.5ex}\titlerule]
\titlespacing*{\section}{0pt}{1.2ex plus .5ex}{0.6ex}
\titlespacing*{\subsection}{0pt}{1ex plus .4ex}{0.4ex}
\titlespacing*{\chapter}{0pt}{-35pt}{14pt}
\titleformat{\section}{\normalfont\Large\bfseries\color{bleuSentinel}}{\thesection}{0.8em}{}
\titleformat{\subsection}{\normalfont\large\bfseries}{\thesubsection}{0.8em}{}
\titleformat{\subsubsection}{\normalfont\normalsize\bfseries\itshape}{\thesubsubsection}{0.8em}{}
\setcounter{secnumdepth}{3}
\setcounter{tocdepth}{1}

% ---- Liens (gris/noir : impression propre) ----------------------------
\usepackage{xurl}
\usepackage[hidelinks,pdfusetitle]{hyperref}
\hypersetup{pdftitle={\titreProjet}, pdfauthor={\etudiant}}

% ---- Abréviations ---------------------------------------------------
\usepackage[nohyperlinks]{acronym}

% ---- Macros du rapport ----------------------------------------------
% Texte à compléter : en rouge entre crochets en mode brouillon (\brouillonfalse pour l'enlever)
\newcommand{\afaire}[1]{\ifbrouillon{\color{rougeAfaire}\bfseries[\,#1\,]}\else #1\fi}

% Capture d'écran : insère l'image si elle existe, sinon un cadre
% indiquant le nom exact du fichier à déposer dans figures/captures/
%   \capture[largeur]{fichier.png}{légende}{label}
\newcommand{\capture}[4][0.95\linewidth]{%
  \begin{figure}[htbp]
    \centering
    \IfFileExists{figures/captures/#2}{%
      \fcolorbox{grisClair!60!black}{white}{\includegraphics[width=#1,height=0.3\textheight,keepaspectratio]{figures/captures/#2}}%
    }{%
      \fcolorbox{rougeAfaire}{rougeAfaire!5}{%
        \begin{minipage}[c][0.45\linewidth][c]{#1}\centering
          \textcolor{rougeAfaire}{\textbf{Capture à insérer}}\\[0.6em]
          \texttt{figures/captures/\detokenize{#2}}\\[0.4em]
          \small\textcolor{grisTexte}{#3}
        \end{minipage}}%
    }
    \caption{#3}\label{#4}
  \end{figure}}

% Diagramme généré (PDF vectoriel dans figures/diagrammes/)
\newcommand{\diagramme}[4][0.95\linewidth]{%
  \begin{figure}[htbp]
    \centering
    \includegraphics[width=#1,keepaspectratio,height=0.36\textheight]{figures/diagrammes/#2}
    \caption{#3}\label{#4}
  \end{figure}}

% Encadré « À retenir »
\newtcolorbox{aretenir}[1][À retenir]{colback=bleuSentinel!5, colframe=bleuSentinel,
  fonttitle=\bfseries, title=#1, boxrule=0.6pt, arc=2pt, left=6pt, right=6pt, top=4pt, bottom=4pt}

% Introduction / conclusion de chapitre
\newenvironment{introchapitre}{\par\noindent\textbf{Introduction}\par\noindent\itshape}{\par\medskip}
\newenvironment{conclusionchapitre}{\par\medskip\noindent\textbf{Conclusion}\par\noindent\itshape}{\par}

% Sous-capture pour une grille de captures (dans une figure)
%   \souscapture[largeur]{fichier.png}{légende}
\newcommand{\souscapture}[3][0.48\linewidth]{%
  \begin{subfigure}[t]{#1}\centering
  \IfFileExists{figures/captures/#2}{\fcolorbox{grisClair!60!black}{white}{\includegraphics[width=0.97\linewidth,height=0.22\textheight,keepaspectratio]{figures/captures/#2}}}{%
  \fcolorbox{rougeAfaire}{rougeAfaire!5}{\begin{minipage}[c][0.52\linewidth][c]{0.93\linewidth}\centering\footnotesize\textcolor{rougeAfaire}{\textbf{Capture à insérer}}\\[0.3em]\texttt{\detokenize{#2}}\end{minipage}}}%
  \caption{#3}\end{subfigure}}

\newcommand{\code}[1]{\texttt{\detokenize{#1}}}
